\section{Implementation details}
\label{app:details}

\paragraph{Training.} We use AdamW with learning rate $10^{-3}$ and weight decay $10^{-4}$, 300 warm-up
steps followed by cosine decay, and gradient clipping at 1.0. Training runs in fp32 with TF32 matrix
multiplies, which was faster than bf16 autocast for this launch-bound model on our GPU. For T1 and T2 the
backbone has width 128 and $r=1$ and is trained for 8000 iterations at batch size 512 with
$\Ttrain=20$. For T3 it has width 64 and $r=4$ and is trained for 5000 iterations at batch size 32, with
random dihedral augmentation and $\Ttrain=32$. Training batches are sampled on the GPU and filtered
against the ID-val and ID-test hashes. Base models use seeds 0--4 and jointly trained methods seeds 0--2
(T1/T2); mazes use seeds 0--1.

\paragraph{Gate hyperparameters.} Every jointly trained gate receives the same budget of three settings
on addition (seed 0). Among the settings within 0.5 points of the best ID-val accuracy, the one with the
fewest mean ID-val steps is used everywhere (\texttt{configs/tune\_gates.yaml}). The grids are
$\tauact\in\{10^{-3},10^{-2},10^{-1}\}$ for the ACT ponder weight,
$(\beta,\lambda_p)\in\{(0.01,0.1),(0.01,0.2),(0.1,0.05)\}$ for PonderNet, and
$\lambda\in\{10^{-3},10^{-2},10^{-1}\}$ for the compute weight of the v1 gate.

\paragraph{Post-hoc heads.} The post-hoc heads are fitted on the frozen backbone for 1500 iterations
(batch 256; batch 32 for mazes) on fresh training-distribution batches rolled out to $2\Ttrain$ steps,
with the backbone in bf16 autocast and the heads in fp32. We use Adam with learning rate
$2\cdot 10^{-3}$ and cosine decay. The pooled correctness head (TRM-style) predicts exact match from the
masked mean of the latent state, and the completion heads predict the target of \cref{eq:done} per
position.

\paragraph{Halting loop.} \Cref{alg:halt} spells out the halting rule of \cref{eq:halt} as it runs at
test time. Instances are processed in batches, and an instance leaves the batch as soon as it halts.

\begin{algorithm}[h]
  \caption{Halting with local completion heads (batched; halted instances leave the batch).}
  \label{alg:halt}
  \begin{algorithmic}[1]
    \Require input $x$, backbone $(E,H,U,D)$, completion head $c$, threshold $\theta$, budget $\Tmax$
    \State $e \gets E(x)$;\quad $h \gets H(e)$;\quad $p \gets \mathbf 0$
    \For{$t = 1,\dots,\Tmax$}
      \State $h \gets U(h, e)$;\quad $\yhat \gets \arg\max D(h)$
      \State $p \gets \sigma\bigl(c(h, p)\bigr)$ \Comment{per position; the pointwise head ignores $p$}
      \If{$\min_{i\ \text{valid}} p_i > \theta$ \textbf{or} $t = \Tmax$}
        \State \Return $(\yhat, t)$
      \EndIf
    \EndFor
  \end{algorithmic}
\end{algorithm}

\paragraph{Threshold grids.} Probability rules use the thresholds
\begin{equation*}
  \theta\in\{0.5,\,0.6,\,0.7,\,0.8,\,0.9,\,0.95,\,0.98,\,0.99,\,0.995,\,0.999,\,0.9999\}.
\end{equation*}
Entropy, KL and hidden-state change
use 19 log-spaced values in $[10^{-6},1]$, stability uses $k\in\{2,3,5\}$, ACT thresholds its cumulative
halting probability at $\{0.5,0.7,0.9,0.95,0.99\}$, PonderNet at $\{0.3,0.5,0.7,0.9,0.95,0.99\}$, and the
fixed-step rule uses $T\in\{1,\dots,\Tmax\}$. The test-time budget of each split is
\begin{equation*}
  \Tmax = \ceil*{\frac{\max_x \dstar(x)}{r}} + 16.
\end{equation*}

\paragraph{Metrics.} Besides accuracy, mean steps and the premature rates of \cref{eq:premature}, we
record several diagnostics. Let $t^{\mathrm{stab}}(x)$ be the first step from which the answer stays
correct, the step chosen by the stable oracle. \emph{Overthinking} is
$\tau(x)-t^{\mathrm{stab}}(x)$, averaged over correctly answered instances. \emph{Regret} is the accuracy
lost and the extra steps spent relative to the stable oracle, the latter measured on instances that the
oracle gets right. \emph{Depth tracking} is the Spearman correlation $\rho_S(\tau,\dstar)$, and
\emph{length leakage} the partial Spearman correlation $\rho_S(\tau, L\mid\dstar)$; a rule that halts at
$\LB(x)$ depends on depth alone and therefore leaks no length. Calibration is the 15-bin expected
calibration error of the per-position done probabilities.

\paragraph{Compute.} \Cref{tab:compute} lists the measured GPU time of each sweep.

\IfFileExists{tables/compute.tex}{\begin{table}[h]
\centering
\caption{Measured GPU time per sweep (RTX 5070; several runs shared the GPU, so wall-clock per run is higher than in isolation). Runs listed in more than one sweep are counted once per sweep.}
\label{tab:compute}
\small
\begin{tabular}{lrrr}
\toprule
Sweep & Runs & Train (h) & Eval (h) \\
\midrule
\texttt{ablations} & 3 & 0.6 & 0.2 \\
\texttt{main\_addition} & 8 & 2.6 & 0.6 \\
\texttt{main\_gates} & 12 & 5.2 & 1.3 \\
\texttt{main\_parity} & 8 & 2.4 & 1.3 \\
\texttt{pilot\_addition} & 8 & 2.4 & 1.1 \\
\texttt{tune\_gates} & 2 & 1.1 & 0.2 \\
\midrule
Total & 41 & 14.3 & 4.8 \\
\bottomrule
\end{tabular}
\end{table}}{}
